\section{The Calibration Pipeline}
\label{sec:pipeline}

The frozen bank is produced by a two-phase pipeline. Phase~1 reads the
generated items, administers them to a grid of synthetic examinees, and
fits Rasch. Phase~2 refits and applies item-fit quality control.

\subsection{Phase 1: the examinee grid}
Phase~1 builds a $3 \times 5 \times 3 = 45$-cell grid of virtual
examinees. The three factors are base model
(\texttt{llama-3.3-70b-versatile}, \texttt{openai/gpt-oss-20b}, and
\texttt{llama-3.1-8b-instant}), sampling temperature
($0.1, 0.4, 0.7, 1.0, 1.2$), and system-prompt strictness
(\texttt{none}, \texttt{neutral}, \texttt{strict}). Every
(examinee~$\times$~item) pair is administered once and graded, which
yields a binary response matrix. On $1{,}823$ items this comes to
$82{,}035$ graded attempts. Each cell is a single zero-shot call, the run
is resumable, and one audit row is written per pair.

\subsection{The dual-mode judge}
Responses are graded by an LLM judge in one of two modes. The
\emph{strict} judge, used for hard and real items, requires the correct
compliant or non-compliant conclusion \emph{and} a section-level citation
such as ``21~CFR~11.10(e)''. The \emph{lenient} judge, used for easy
items, requires only the correct conclusion and consistent reasoning,
with no citation. The lenient mode is necessary because without it the
easy items would be failed by weak examinees that never cite sections,
which would collapse the low-difficulty anchor. All judge errors are
mapped to a failing score, so that no exception escapes the matrix.

\subsection{Rasch fit and baseline anchoring}
I fit a 1PL Rasch model by marginal maximum likelihood
(\texttt{girth.rasch\_mml}; \citealp{girth}), with discrimination fixed
at $a = 1$ for all items. Difficulty is identified only up to an additive constant, so the
scale has to be anchored. In the 1PL model
$P(\text{correct}\mid\theta,b)=\sigma(\theta-b)$, so shifting every $b$
by a constant $c$ shifts every $\theta$ by $-c$. I choose
$c = \hat{\theta}_{\text{base}}$, the MLE ability of a designated
baseline examinee (\texttt{llama-3.1-8b-instant}, $t{=}0.1$, strictness
\texttt{none}). The baseline then sits at $\theta = 0$ by construction,
and all reported $b$ values are on that scale.

\subsection{Phase 2: classical item-fit quality control}
Phase~2 refits Rasch on the Phase~1 matrix and drops any item that fails
one of four classical checks.

\begin{enumerate}[leftmargin=1.4em,topsep=2pt]
  \item \textbf{Zero variance.} All-pass or all-fail items carry no
  information and their difficulty is undefined at the bound. These are
  dropped unconditionally.
  \item \textbf{Point-biserial} (corrected, item versus rest-of-total).
  A non-positive correlation means that stronger examinees do
  \emph{worse} on the item, which signals a broken gold standard. These
  are dropped at $\le 0$.
  \item \textbf{Infit} (information-weighted mean-square residual),
  dropped above $1.5$.
  \item \textbf{Outfit} (unweighted mean-square residual), dropped above
  $2.0$.
\end{enumerate}

The fit thresholds are deliberately asymmetric, with upper bounds only.
Linacre's conventional bands \citep{linacre2002infit} treat mean-squares
in $[0.5, 1.5]$ as productive for measurement, $1.5$ to $2.0$ as
unproductive but not degrading, and above $2.0$ as degrading. That
symmetric guidance is calibrated for thousands of examinees. On a 45-cell
grid, an item with a very low mean-square is a clean separator that I
want to keep, not a defect. Outfit gets the looser bound ($2.0$ against
infit's $1.5$) because it is unweighted, so a single anomalous cell at
extreme $\theta$ inflates it disproportionately on a small grid.

\paragraph{Why Rasch and not 2PL.}
A 2PL model is not identifiable on a 45-examinee grid. Marginal MLE
pinned every discrimination at the optimizer's upper bound, and joint
MLE, although it gives a real spread, is statistically inconsistent as
the number of items grows with the number of examinees held fixed. That
is exactly the regime here. Rasch is identifiable, it gives the specific
objectivity that the $\Delta\theta$ comparisons in
Section~\ref{sec:deltatheta} rely on, and it matches the model used
during generation. I retain the 2PL machinery for a future
re-evaluation, once the examinee count grows.
